% SPDX-FileCopyrightText: 2026 Will Cook
% SPDX-License-Identifier: CC-BY-4.0
%
% Synthesis note: mathematics whose subject is several of the Erdős Problem
% Notes read together.  Build with `make -C paper` or
% `tectonic erdos-synthesis-subsums-across-bases.tex`.
\documentclass[11pt]{article}
% SPDX-FileCopyrightText: 2026 Will Cook
% SPDX-License-Identifier: CC-BY-4.0
%
% Shared preamble for the Erdős Problem Notes: the short, standalone,
% problem-owned manuscripts covering the expansion library `ErdosProblems`.
%
% One note owns one Erdős problem.  Each note repeats whatever context its own
% reader needs, so that a reader who arrives at a single problem never has to
% assemble it from the gateway paper.  Deliberate repetition across notes is the
% design, not an oversight.
%
% Source links resolve against one pinned formal-source commit, declared once
% below.  `scripts/check_problem_note_sources.py` verifies that every authored
% (file, line, declaration) triple names that exact declaration in the pinned
% snapshot, so a link cannot silently rot as later work moves lines.

\usepackage{paper-house-style}
\usepackage{array}
\usepackage{booktabs}
\usepackage{enumitem}
\setlist{topsep=0.35em,itemsep=0.2em}
\usepackage[colorlinks=true,linkcolor=linkinternal,urlcolor=linkexternal,
  citecolor=linkinternal]{hyperref}
\hypersetup{bookmarksnumbered=true,bookmarksopen=true,bookmarksopenlevel=1,
  pdfauthor={Will Cook},pdflang={en-GB}}

% ---------------------------------------------------------------------------
% Pinned formal source.  \commit names the checkpoint every link resolves
% against; the checker reads that snapshot rather than the working tree, so the
% printed coordinates stay correct however later waves move declarations.
% ---------------------------------------------------------------------------
\newcommand{\commit}{e5fd26a6d1b1a346430205eab5385b4c9565d004}
\newcommand{\commitshort}{e5fd26a6d1b1}
\newcommand{\repobase}{https://github.com/wcook04/plectis-lean-erdos249-257/blob/\commit}
\newcommand{\sourceurl}{https://github.com/wcook04/plectis-lean-erdos249-257/tree/\commit}
\newcommand{\PX}{ErdosProblems}

% Declaration names stay inside link targets.  The printed label is either a
% spaced, lower-case reading of the declaration name (\lref) or a short authored
% phrase (\lword); neither prints a module path or a raw Lean identifier.
\ExplSyntaxOn
\cs_new_protected:Npn \noteleanlabel #1
  {
    \tl_set:Nx \l_tmpa_tl { \tl_to_str:n {#1} }
    \regex_replace_all:nnN { _+ } { \c{space} } \l_tmpa_tl
    \regex_replace_all:nnN
      { ([a-z0-9])([A-Z]) }
      { \1\c{space}\2 }
      \l_tmpa_tl
    \tl_set:Nx \l_tmpa_tl { \tl_lower_case:n { \tl_use:N \l_tmpa_tl } }
    \tl_use:N \l_tmpa_tl
  }
\ExplSyntaxOff
\newcommand{\leanlabel}[1]{\noteleanlabel{#1}}
\newcommand{\lref}[3]{\href{\repobase/\PX/#1\#L#2}{\leanlabel{#3}}}
\newcommand{\lword}[4]{\href{\repobase/\PX/#1\#L#2}{#4}}
\newcommand{\lloc}[2]{\href{\repobase/\PX/#1\#L#2}{Lean source}}

% Links into a sibling library, named repository-relative.  The expansion
% library is implicit in \lref/\lword, because that is where problem-owned
% work lands.  Several problems have their headline theorem in the older
% reviewed #249/#257 corpus, and a note that could not name that library
% would have to paraphrase a checked statement instead of linking it.
\newcommand{\mref}[3]{\href{\repobase/#1\#L#2}{\leanlabel{#3}}}
\newcommand{\mword}[4]{\href{\repobase/#1\#L#2}{#4}}
\newcommand{\mloc}[2]{\href{\repobase/#1\#L#2}{Lean source}}
\emergencystretch=2em

\newtheorem{theorem}{Theorem}[section]
\newtheorem{proposition}[theorem]{Proposition}
\newtheorem{lemma}[theorem]{Lemma}
\newtheorem{corollary}[theorem]{Corollary}
\theoremstyle{definition}
\newtheorem{definition}[theorem]{Definition}
\newtheorem{example}[theorem]{Example}
\newtheorem{problem}[theorem]{Problem}
\theoremstyle{remark}
\newtheorem*{remark}{Remark}

\newcommand{\N}{\mathbb{N}}
\newcommand{\Npos}{\mathbb{N}_{>0}}
\newcommand{\Z}{\mathbb{Z}}
\newcommand{\Q}{\mathbb{Q}}
\newcommand{\R}{\mathbb{R}}
\newcommand{\lcm}{\operatorname{lcm}}
\newcommand{\ph}{\varphi}

% The series line printed under every note's title block.
\newcommand{\seriesline}{Erd\H{o}s Problem Notes}

% Production provenance.  The wording is fixed by the corpus provenance
% contract and reproduced verbatim in every manuscript in this repository, so
% the papers cannot drift into different accounts of how they were made.
\newcommand{\noteprovenance}{\provenancenote{\emph{Authorship and AI use.} The
prose, formal proofs, and software in this version were drafted and revised by
large-language-model agents under Will Cook's direction.  Cook set the
objectives and acceptance criteria, reviewed and authorised the manuscript, and
is responsible for its contents.  The agents are production tools, not authors.
See \emph{Statements and declarations}.}}

% The status paragraph every note carries, in its own words, before any result.
% Kernel checking establishes the exact Lean proposition; it does not establish
% that the proposition is the interesting one, that it is new, or that the
% Erdős problem is closed.
\newcommand{\statusboundary}{%
  \emph{Status.}  The problem treated here is open, and this note does not
  close it.  Every statement below marked as checked is a proposition that the
  pinned Lean kernel accepts from the sources this note links to, with no
  \texttt{sorry}, no added axiom, and no unchecked evaluation.  That is a claim
  about the formal statement, not about its mathematical interest, its
  novelty, or the original problem.  The unresolved obligations are named
  exactly, in their own section, and none of the finite computations,
  reductions, or no-go results here removes one of them.}

\usepackage{xurl}
\hypersetup{
  pdftitle={Rational and irrational subsums of a Lambert series across bases},
  pdfsubject={Erdős Problem Notes: synthesis across problems 249, 251, 257 and 1049},
  pdfkeywords={irrationality, Lambert series, subsum sets, achievement sets, Mahler's method, Farey fractions},
}

% Immutable formal-source revision used by the corpus; this note links no Lean
% source itself and cites sibling papers by file name.
\renewcommand{\commit}{99f4bf47422abbd8757cbb22b50ba079d764d3a7}
\renewcommand{\commitshort}{99f4bf47422a}

\runninghead{Subsums of a Lambert series across bases}
\title{Rational and Irrational Subsums\\of a Lambert Series Across Bases}
\subtitle{Choices against contraction in seven irrationality problems of Erd\H{o}s}
\author{Will Cook}
\date{20 September 2026}

\begin{document}
\maketitle
\noteprovenance

\begin{abstract}
Let $a>b\ge1$ be coprime integers with $a^2>b^3$.  We prove that
$\sum_{n\in S}((a/b)^n-1)^{-1}$ is irrational for every infinite set $S$ of
positive integers that is totally ordered by divisibility, and transcendental
at every rational base when $S$ eventually doubles.  For every rational base
strictly between $1$ and $2$ the same sums take rational values on infinite
sets $S$; at base $3/2$ the value $1/2$ occurs.  The interval result follows
from subsum covering and a denominator observation; the chain result combines
nested denominators with Mahler's method.  Subsum sets are null when the choices fall short
by an unbounded factor and contain intervals when every step is at most the
reach of the later steps, and base $2$ lies exactly on the boundary.  On the
interval side no property shared by all the competing series can decide
irrationality; the rational series constructed in the companion notes on
Erd\H{o}s Problems \#249 and \#251 are instances.  At base $2$ the subsums form
a Cantor set of measure one.  We determine the limiting proportion of rationals
of bounded height not rejected through any fixed number of greedy steps.
This proportion does not settle membership; an exact late rejection shows
why it gives no deterministic stopping depth for the computation.
\end{abstract}

\section{Results}\label{sec:results}

Seven of the eight problems treated in the companion notes ask whether a
series is irrational: $\sum(n!-1)^{-1}$ in \#68, reciprocal sums of
near-Sylvester sequences in \#243, $\sum\varphi(n)2^{-n}$ in \#249,
$\sum p_n2^{-n}$ in \#251, the subseries $\sum_{n\in S}(2^n-1)^{-1}$ in \#257,
a running least common multiple in \#269, and $\sum(t^n-1)^{-1}$ at rational
$t$ in \#1049 \cite{erdosgraham1980,erdos1988}.  The eighth, \#1041, is
geometric and plays no part here.  This note reads the seven together.  Its
theorems concern the two that share a series.  For real $t>1$ and a set $S$ of
positive integers put
\[
 X_S(t)=\sum_{n\in S}\frac{1}{t^n-1},\qquad
 w_n=\frac{1}{t^n-1},\qquad R_N=\sum_{n>N}w_n .
\]
Problem \#257 asks whether $X_S(t)$ is irrational for every infinite $S$ at
every integer $t\ge2$ \cite[p.~62]{erdosgraham1980}; Problem \#1049 asks about
$X_{\N_{>0}}(t)$ at rational $t$.

The comparison behind everything is stated for weights with multiplicities.
Let $u_n>0$ and integers $D_n\ge1$ satisfy $\sum_nD_nu_n<\infty$, and put
\[
 V=\Bigl\{\sum_{n\ge1}\varepsilon_nu_n:\ \varepsilon_n\in\{0,1,\dots,D_n\}\Bigr\},
 \qquad C_N=\sum_{n>N}D_nu_n .
\]

\begin{theorem}[choices against contraction]\label{thm:dichotomy}
\textup{(i)} For every $N$, the Lebesgue measure of $V$ is at most
$C_N\prod_{n\le N}(D_n+1)$.  If $\liminf_NC_N\prod_{n\le N}(D_n+1)=0$ then $V$
is null.

\textup{(ii)} If $u_n\le C_n$ for every $n>N_0$, then for each choice of
$\varepsilon_1,\dots,\varepsilon_{N_0}$ the set $V$ contains the interval
$[\mu,\mu+C_{N_0}]$, where $\mu=\sum_{n\le N_0}\varepsilon_nu_n$.

\textup{(iii)} For $u_n=\beta^{-n}$ with real $\beta>1$ and $D_n=D$, the bound in
\textup{(i)} tends to $0$ exactly when $D+1<\beta$, and the hypothesis of
\textup{(ii)} holds exactly when $D+1\ge\beta$.
\end{theorem}

Part (ii) is Kakeya's covering argument and part (i) is the standard covering
bound; both are classical \cite{hornich1941,nitecki2013,fridy1966}, and
Kova\v{c} and Tao give a scalar reciprocal-choice covering lemma
\cite[Lemma~5.1]{kovactao2024}; their higher-dimensional approximation lemma is
Lemma~7.2 of the same paper.  We claim no novelty for
Theorem~\ref{thm:dichotomy}.  Its use here is to say which side each problem
lies on.

\begin{theorem}[the subseries across bases]\label{thm:bases}
Let $t>1$ be real.

\textup{(a)} If $t<2$, let $N_0\ge0$ be least with $t^{-n}\le2-t$ for all
$n>N_0$.  Then the set of values $X_S(t)$ is the union of the intervals
$[X_F(t),X_F(t)+R_{N_0}]$ over $F\subseteq\{1,\dots,N_0\}$, where
$X_F(t)=\sum_{n\in F}w_n$; in particular it contains $[0,R_{N_0}]$.  If moreover
$t=a/b$ in lowest terms, then every rational in $[0,R_{N_0}]$ whose reduced
denominator shares a prime factor with $ab$ equals $X_S(t)$ for some $S$, and
every such $S$ is infinite.  There is an infinite $S\subseteq\{2,3,\dots\}$
with $X_S(3/2)=1/2$.

\textup{(b)} If $t=2$, every value has exactly one $S$, and the set of values
is a Cantor set of Lebesgue measure $1$ inside $[0,E]$, where
$E=\sum_{n\ge1}(2^n-1)^{-1}=1.6066951524\ldots$.

\textup{(c)} If $t>2$, the set of values is null.
\end{theorem}

Part (b) is proved in the companion note on \#257, Section~7, from Hornich's
theorem as proved by Nitecki \cite{plectis257,hornich1941,nitecki2013}; Kova\v{c}
and Tao record the strict inequality $w_N>R_N$ and the Cantor-set conclusion for
every fixed base $t\ge2$ \cite[Remark~4.1]{kovactao2024}.  Parts (a) and (c) follow
from Theorem~\ref{thm:dichotomy} in a few lines.  We have not found part (a)
stated for non-integer bases and it may be known.  It shows that the statement
asked in \#257 is false at every rational base below $2$, so any proof at base
$2$ must use more than the shape of the series.

\begin{theorem}[divisibility chains]\label{thm:chains}
Let $t=a/b>1$ in lowest terms and let $S=\{n_1<n_2<\cdots\}$ be infinite with
$n_j\mid n_{j+1}$ for every $j$.

\textup{(a)} If $a^2>b^3$, then $X_S(t)$ is irrational.

\textup{(b)} If $n_{j+1}=2n_j$ for all large $j$, then $X_S(t)$ is
transcendental, with no condition on $a$ and $b$.
\end{theorem}

The hypothesis $a^2>b^3$ says $\log b/\log a<2/3$.  It holds for every integer
base, for $3/2$, $5/2$ and $7/3$, and fails for $4/3$ and $5/4$.  At integer
bases part (a) is contained in the theorem of Erd\H{o}s on supports with
$\sum_{n\in S}1/n<\infty$ \cite{erdos1968}, proved in full in
\cite[Theorem~2]{plectis257}.  At base $3/2$ it sits inside the regime of
Theorem~\ref{thm:bases}(a): rational values occur there, and exact divisibility
still forces irrationality.  For comparison, the companion note on \#1049 proves the irrationality of the
full sum $X_{\N_{>0}}(a/b)$ when $\log b/\log a<0.4056830213840605\ldots$, using Zudilin's
linear forms \cite[Theorem~1]{plectis1049}, \cite{zudilin2004}, and proves that
the sufficient cutoff supplied by one integer-polynomial family with common
leading degree, coefficient-height and decay bounds at every fixed real base
$x>1$ is at most $1/2$ \cite[Theorem~5]{plectis1049}.  This restriction does not
exclude stronger estimates at a particular base or a different choice of
family there.  Thin supports reach further than the full sum
because the denominators divide one another.

At base $2$ the greedy rule characterises membership and supplies finite
certificates of nonmembership: starting from $r_0=x$, take
index $n$ when $r_{n-1}\ge w_n$ and subtract.  Since $w_n>R_n$, a real
$x\in[0,E]$ is a subsum if and only if no remainder falls strictly between
$R_n$ and $w_n$; we say $x$ is \emph{rejected at step $n$} when that happens
first at index $n$.

\begin{theorem}[the fixed-depth rational count]\label{thm:forced}
Fix $N\ge1$.  Among the reduced fractions $p/q\in(0,E]$ with $q\le Q$, the
proportion not rejected in the first $N$ steps tends to $2^NR_N/E$ as
$Q\to\infty$, with error $O(2^N\log Q/Q)$.  Consequently the upper limit of
the proportion that are subsums is at most $1/E=0.62239\ldots$.
\end{theorem}

The limit $2^NR_N/E$ does not depend on whether \#257 is true.  Agreement with
this fixed-depth limiting proportion therefore does not establish membership.
An exact computation in \cite{plectisinvestigation} first read a surviving
share near $62\%$ as evidence that most such fractions are subsums;
Theorem~\ref{thm:forced} is the correction.  Section~\ref{sec:base2} states
what remains.

\paragraph{Evidence.}  Every proof in this note is an ordinary proof.  Nothing
here is checked in Lean, and nothing has been reviewed by a specialist.
Statements quoted from the companion notes carry the status given there, which
we repeat at each use.  Theorem~\ref{thm:chains} uses Nishioka's theorem for Mahler systems as an
external premise in the eventually doubling case, hence in both parts.

\section{Choices against contraction}\label{sec:dichotomy}

\begin{proof}[Proof of Theorem~\ref{thm:dichotomy}]
(i) A choice of $\varepsilon_1,\dots,\varepsilon_N$ fixes
$\sum_{n\le N}\varepsilon_nu_n$, and the rest of the sum lies in $[0,C_N]$.
So $V$ is covered by at most $\prod_{n\le N}(D_n+1)$ intervals of length $C_N$.

(ii) Subtracting $\mu$, it suffices to reach every $y\in[0,C_{N_0}]$ with the
levels $n>N_0$.  Put $y_{N_0}=y$ and, for $n>N_0$, let
$\varepsilon_n=\min(D_n,\lfloor y_{n-1}/u_n\rfloor)$ and
$y_n=y_{n-1}-\varepsilon_nu_n$.  If $0\le y_{n-1}\le C_{n-1}=D_nu_n+C_n$ then
$0\le y_n\le C_n$: when $\varepsilon_n=D_n$ this is a subtraction, and when
$\varepsilon_n<D_n$ it holds because $y_n<u_n\le C_n$.  Since $C_n\to0$, the
sum of the $\varepsilon_nu_n$ is $y$.

(iii) Here $C_n=D\beta^{-n}/(\beta-1)$.  The product in (i) is a constant
multiple of $((D+1)/\beta)^N$, and $u_n\le C_n$ says $\beta-1\le D$.
\end{proof}

The reading that matters for irrationality is immediate.  Suppose a family of
series has values $x_0+\sum\varepsilon_nu_n$ with the $\varepsilon_n$ free as in
(ii).  The values then fill an interval, which contains rationals and
irrationals, so a property that all members of the family share implies
neither.  A proof of irrationality for one member has to use something that
distinguishes it inside the family.  Two constructions in the companion notes
are families of this kind.

\emph{Problem \#251.}  Proposition~1 of \cite{plectis251} (ordinary proof),
applied in its Corollary~2 to the prime gaps, changes the gaps on a set of upper
Banach density zero, by nonnegative integers below any prescribed function
tending to infinity, keeping every congruence modulo every fixed $q$ from some
point on, and reaches every real in an interval.  The resulting positions
$P_n$ satisfy $P_n\sim n\log n$ and are not asserted to be prime.  In the notation above the $j$-th block has $u_j=M_j2^{-n_j-2}$ and
$D_j=2^{s_j}-1$, and the inequality verified there is $u_j\le C_j$.  The
lemma is credited there to \cite{fridy1966}, \cite[Lemma~4]{crmarickovac2025}
and \cite[Lemma~5.1]{kovactao2024}.  The note also proves that a bounded
allowance is impossible under the stated congruences.  Thus every allowance
tending to infinity suffices, while no bounded allowance does.  The specified
growth, eventual fixed-modulus congruences and empirical distributions of
unnormalised blocks therefore do not suffice to prove
$\sum p_n2^{-n}$ irrational; neither primality nor every quantitative
correlation is preserved.

\emph{Problem \#249.}  Section~5 of \cite{plectis249} gives an integer sequence
$c$ with $c(n)=\varphi(n)$ for odd $n$, $|c(n)-\varphi(n)|\le2$ for even $n$,
$0\le c(n)\le n$ and $\sum c(n)2^{-n}=5/4$ (Lean-checked there).  Changing even
indices by at most $2$ is the case $u_n=2^{-n}$ for even $n$, $D_n=4$, where
$C_n\ge\tfrac43 2^{-n}>u_n$.

The same comparison appears in Kova\v{c} and Tao's theorem that for integers
$2\le t_1<\dots<t_m$ with $\sum1/(t_k-1)>1$ there are sets $S_k$, one of them
infinite, with $\sum_kX_{S_k}(t_k)$ rational \cite[Theorem~2.3]{kovactao2024}:
merging the weights of the $m$ bases, the sum of all weights below a given
weight $u$ is at least $(1-O(u))\,u\sum1/(t_k-1)$, which exceeds $u$ once $u$
is small.

\section{The subseries across bases}\label{sec:bases}

\begin{proof}[Proof of Theorem~\ref{thm:bases}]
(a) Since $R_n\ge\sum_{k>n}t^{-k}=t^{-n}/(t-1)$ and
$w_n=t^{-n}/(1-t^{-n})$, the inequality $w_n\le R_n$ holds as soon as
$t^{-n}\le2-t$.  Theorem~\ref{thm:dichotomy}(ii) with $D_n=1$ gives, for
each $F\subseteq\{1,\dots,N_0\}$, the interval $[X_F(t),X_F(t)+R_{N_0}]$, and every
value $X_S(t)$ lies in the one with $F=S\cap[1,N_0]$.  For $t=a/b$ a finite subsum is
$\sum_{n\in F}b^n/(a^n-b^n)$, whose denominator divides $\prod(a^n-b^n)$ and is
coprime to $ab$.  A rational whose denominator is not coprime to $ab$ is
therefore never a finite subsum, and inside the interval it is a subsum.  At
$t=3/2$ the inequality $t^{-n}\le1/2$ holds for $n\ge2$ and $R_1>1/2$.

(c) For $t>2$, $2^NR_N\le2^Nt^{-N}/((t-1)(1-t^{-1}))\to0$, and
Theorem~\ref{thm:dichotomy}(i) applies.
\end{proof}

\begin{proof}[Proof of Theorem~\ref{thm:chains}]
Write $\rho=b/a$, $D_j=a^{n_j}-b^{n_j}$ and
$S_j=\sum_{i\le j}b^{n_i}/(a^{n_i}-b^{n_i})$.  Because $n_i\mid n_j$ for
$i\le j$, each $a^{n_i}-b^{n_i}$ divides $D_j$, so $D_jS_j$ is an integer.  The
tail satisfies
\[
 0<X_S(t)-S_j\le\sum_{n\ge n_{j+1}}\frac{\rho^n}{1-\rho^n}
 \le\frac{\rho^{\,n_{j+1}}}{(1-\rho)^2}.
\]
Each ratio $n_{j+1}/n_j$ is an integer at least $2$.

Suppose the ratio is at least $3$ for infinitely many $j$.  For those $j$,
\[
 0<D_j\bigl(X_S(t)-S_j\bigr)\le\frac{a^{n_j}\rho^{3n_j}}{(1-\rho)^2}
 =\frac{1}{(1-\rho)^2}\Bigl(\frac{b^3}{a^2}\Bigr)^{n_j}\longrightarrow0
\]
when $a^2>b^3$.  If $X_S(t)=p/q$, then $qD_j(X_S(t)-S_j)$ is a positive
integer for every $j$, a contradiction.

Otherwise the ratio is $2$ from some index on, which is the case of part (b),
and $S$ is a finite set together with $\{d,2d,4d,\dots\}$.  With $z=t^{-d}$ the
infinite part contributes $g(z)$, where
\[
 g(z)=\sum_{k\ge0}\frac{z^{2^k}}{1-z^{2^k}},\qquad
 g(z^2)=g(z)-\frac{z}{1-z}.
\]
We use Nishioka's value theorem for Mahler systems
\cite{nishioka1996}, in the precise form quoted by Adamczewski and Faverjon
\cite[Theorem~1.1, p.~3]{adamczewskifaverjon2016}.  At an algebraic regular
point it equates the transcendence degree of the function values with that
of the functions over $\overline{\Q}(z)$.  Here the system is
\[
 \begin{pmatrix}g(z)\\1\end{pmatrix}
 =\begin{pmatrix}1&z/(1-z)\\0&1\end{pmatrix}
  \begin{pmatrix}g(z^2)\\1\end{pmatrix}.
\]
The matrix and its inverse have no poles at any iterate
$\rho^{\,d2^k}\in(0,1)$, so $\rho^d$ is regular.
The function $g$ is transcendental: at a root of unity
$\zeta$ of order $2^j$ the terms with $k\ge j$ are positive on the ray
$r\zeta$ and diverge as $r\to1$, while the earlier terms stay bounded, so $g$
has infinitely many singularities on the unit circle.  An algebraic function
over $\mathbb{C}(z)$ has only finitely many singular points, so $g$ is
transcendental over $\mathbb{C}(z)$.  Hence $g(\rho^{\,d})$ is
transcendental, and adding the rational finite part proves (b) and completes
(a).
\end{proof}

We do not know whether $a^2>b^3$ is necessary in (a).  The first open case is
base $4/3$ with a chain whose ratios are $2$ except for infinitely many $3$s.

\section{Base two}\label{sec:base2}

\begin{proof}[Proof of Theorem~\ref{thm:forced}]
Fixing the first $N$ digits gives $2^N$ closed intervals of length $R_N$.
They are pairwise disjoint because $w_n>R_n$, and a real in $[0,E]$ is not
rejected in the first $N$ steps exactly when it lies in their union $K_N$,
which has measure $2^NR_N$.  For an interval $I$ of length $\ell$, the number
of integers $p$ coprime to $q$ with $p/q\in I$ is
$\varphi(q)\ell+O(2^{\omega(q)})$ by inclusion and exclusion.  Summing over
$q\le Q$ with $\sum_{q\le Q}\varphi(q)=3Q^2/\pi^2+O(Q\log Q)$ and
$\sum_{q\le Q}2^{\omega(q)}=O(Q\log Q)$ gives $3Q^2\ell/\pi^2+O(Q\log Q)$.
Apply this to the $2^N$ intervals of $K_N$ and to $(0,E]$ and divide.  The
second statement follows because every subsum lies in every $K_N$ and
$2^NR_N\to1$.
\end{proof}

The intervals removed at step $n$ are explicit.  Put $g_n=w_n-R_n$, so that
$g_n=\tfrac23 4^{-n}+O(8^{-n})$.  For a finite nonempty $F$ with largest element
$n$, the interval $(X_F(2)-g_n,\,X_F(2))$ is removed at step $n$, and
\[
 [0,E]\smallsetminus\mathcal A=\bigsqcup_{F\ne\varnothing}\bigl(X_F(2)-g_{\max F},\,X_F(2)\bigr),
 \qquad\sum_{n\ge1}2^{n-1}g_n=E-1,
\]
where $\mathcal A$ is the set of subsums.  By Theorem~\ref{thm:bases}(b) a
rational with an infinite $S$ is exactly a counterexample to \#257 at base $2$.
So \#257 at base $2$ holds if and only if every rational in $[0,E]$ that is not
a finite subsum lies strictly between $X_F(2)-g_{\max F}$ and $X_F(2)$ for some
finite nonempty $F$.
This is a one-sided question of approximation by the countable set of finite
subsums, with an error that shrinks like $4^{-\max F}$.

Under \#257 the only fractions of height at most $Q$ that are subsums are the
finite subsums, $40$ of the $19{,}653$ fractions with $2\le q\le200$.  A model
that treats later remainders as equidistributed gives the opposite extreme, a
proportion tending to $1/E$, because the shares $2^{n-1}g_n/E$ of the removed
intervals are summable.  That model asserts that \#257 fails for a
positive proportion of all rationals.  The fixed-depth limiting proportion does
not distinguish this claim from its negation.  An arithmetic argument, or a
count with separately justified estimates as both depth and height grow, is needed.
Measure does not decide either.  Boes, Darst and Erd\H{o}s construct symmetric
Cantor sets of every measure in $[0,1)$ that contain essentially no rationals
\cite{boesdarsterdos1981}.

The exact computation in \cite{plectisinvestigation} agrees with
Theorem~\ref{thm:forced} step by step.  Among the $19{,}653$ reduced fractions
with $2\le q\le200$, the numbers rejected at steps $1$, $2$ and $7$ are $4809$,
$1470$ and $32$, against $4811$, $1467$ and $32$ from the measures of the
removed intervals.  The measure-based main term for the number rejected at step $n$ is
$2^{n-1}g_n\sum_{2\le q\le Q}\varphi(q)$, asymptotically $(3Q^2/\pi^2)2^{n-1}g_n$,
which falls below $1$ near $n=2\log_2Q-2\log_2\pi$.  This is not a
deterministic cutoff: the error in Theorem~\ref{thm:forced} does not justify
such an extrapolation.  For example, $189/388$ is first rejected at step $17$,
beyond this scale for $Q=388$.  Its selected indices before rejection are
$F=\{2,3,7,9,10,14,15,16\}$, and exact arithmetic gives
\[
 R_{17}\le\frac{196609}{25769803776}
 <\frac{189}{388}-X_F(2)
 =\frac{9291822600689}{1217890317075045460}
 <\frac1{131071}=w_{17}.
\]
The working record supplies the earlier-step checks.  Deeper computation can
therefore produce new exclusion certificates; absence of a later rejection
still does not prove membership.  No rational is known to have an infinite
$S$ at base $2$.

\begin{problem}\label{prob:membership}
Decide whether $1/2$ is a subsum of $\sum(2^n-1)^{-1}$.  By
\cite[Theorem~7]{plectis257} this holds if and only if the integer remainders
of the greedy rule fail to increase at infinitely many steps.
\end{problem}

The rational-point counting papers examined here concern null Cantor sets
such as the middle-third set \cite{rstw2019,chowvarjuyu2024}.  We did not locate
a theorem settling the present positive-measure subsum problem.
Problem~\ref{prob:membership} asks about one explicit rational point.

\section{Limits on methods}\label{sec:second}

The companion notes prove a second kind of negative result.  Each bounds what
one method can do, and none follows from Theorem~\ref{thm:dichotomy}, which
concerns sets of values.
\begin{itemize}
\item \#1049, Theorem~5 \cite{plectis1049} (ordinary proof there, with the
  contradiction step and the comparison with $1/2$ checked in Lean).  For one
  family of nonzero integer-polynomial linear forms, assume degrees at most
  $\delta n^2(1+o(1))$, logarithmic coefficient heights at most
  $h n^2(1+o(1))$, and logarithmic decay
  $-\sigma n^2\log x\,(1+o(1))$ at every fixed real $x>1$, with common constants
  $\delta,\sigma>0$ and $h\ge0$.  Then $\sigma\le\delta$, so the sufficient
  region supplied by these estimates,
  $\log b/\log a<\sigma/(\sigma+\delta)$, never reaches $1/2$.
\item \#269, Theorem~1 \cite{plectis269} (Lean-checked there; Fan posted the
  two-prime separation \cite{fan2026comment}, and the three-prime statement is
  proved in the note).  For three primes the reciprocal
  running-LCM kernel has nonsingular minors of every order, so no finite sum of
  products separates one exponent from the other two.
\item \#257, Section~3 \cite{plectis257} (ordinary proof).  Every positive
  divisor cover costs at least $e$ times the mean of $\log^+$ of its
  multiplicity, and the averaging method cannot reach the prime support, where
  irrationality is known at base $2$ by other means \cite{taoteravainen2025}.
\item \#249, Theorem~9 \cite{plectis249} (Lean-checked there).  Every
  admissible rank-one quotient stays at distance more than $21/320$ from its
  target.
\item \#68, Section~6 \cite{plectis68} (ordinary proof).  Under the displayed
  cancellation hypotheses, the integer-gap comparison fails at cutoffs
  $N=D+O(1)$ as the cancellation cutoff $D\to\infty$.
\item \#243, Proposition~17 \cite{plectis243} (Lean-checked there).  A
  counterexample must have relative errors tending to $0$ with unbounded
  negative parts, so size conditions on the error do not decide.
\end{itemize}
Two of these are inequalities between a cost and a gain: $\sigma\le\delta$ in
\#1049 and the cover cost in \#257.  Whether one inequality governs both is
open, and we know of no formulation that covers the other four.

\begin{thebibliography}{99}
\bibitem{erdosgraham1980} P.~Erd\H{o}s and R.~L. Graham,
  \href{https://mathweb.ucsd.edu/~ronspubs/80_11_number_theory.pdf}{\emph{Old
  and New Problems and Results in Combinatorial Number Theory}},
  Monogr. Enseign. Math. 28, Geneva, 1980.

\bibitem{erdos1988} P.~Erd\H{o}s, \emph{On the irrationality of certain series:
  problems and results}, in A.~Baker (ed.), \emph{New Advances in Transcendence
  Theory}, Cambridge UP, 1988, pp.~102--109,
  doi:\href{https://doi.org/10.1017/CBO9780511897184.009}{10.1017/CBO9780511897184.009}.

\bibitem{erdos1968} P.~Erd\H{o}s,
  \href{https://users.renyi.hu/~p_erdos/1969-09.pdf}{\emph{On the
  irrationality of certain series}}, Math.\ Student 36 (1968), 222--226
  (issued 1969).

\bibitem{hornich1941}
H.~Hornich,
\emph{\"Uber beliebige Teilsummen absolut konvergenter Reihen},
Monatshefte f\"ur Mathematik und Physik \textbf{49} (1941), 316--320.
\href{https://doi.org/10.1007/BF01707309}{doi:10.1007/BF01707309}.

\bibitem{nitecki2013} Z.~Nitecki,
  \href{https://arxiv.org/abs/1106.3779v2}{\emph{Subsum sets: intervals, Cantor
  sets, and Cantorvals}}, arXiv:1106.3779v2 (2013).

\bibitem{fridy1966} John A. Fridy, \href{https://www.fq.math.ca/Scanned/4-3/fridy.pdf}{\emph{Generalized bases for the real numbers}}.
  The Fibonacci Quarterly \textbf{4} (1966), no.~3, 193--201.

\bibitem{crmarickovac2025} T.~Crmari\'{c} and V.~Kova\v{c},
  \href{https://arxiv.org/abs/2504.18712v1}
  {\emph{On the irrationality of certain super-polynomially decaying series}},
  Colloq. Math. \textbf{179} (2025), 55--68,
  doi:\href{https://doi.org/10.4064/cm9628-5-2025}{10.4064/cm9628-5-2025}.
  Theorem locators follow arXiv:2504.18712v1.

\bibitem{kovactao2024} Vjekoslav Kova\v{c} and Terence Tao, \href{https://arxiv.org/abs/2406.17593v4}{\emph{On several irrationality problems for Ahmes series}}.
  Acta Mathematica Hungarica \textbf{175} (2025), 572--608, doi:\href{https://doi.org/10.1007/s10474-025-01528-0}{10.1007/s10474-025-01528-0}; arXiv:\href{https://arxiv.org/abs/2406.17593v4}{2406.17593v4}. Statement and section numbers refer to arXiv version 4.

\bibitem{zudilin2004} W.~Zudilin,
  \href{https://geodesic.mathdoc.fr/articles/10.4064/aa111-2-4/}{\emph{Heine's
  basic transform and a permutation group for \(q\)-harmonic series}},
  Acta Arith. \textbf{111} (2004), no.~2, 153--164,
  doi:\href{https://doi.org/10.4064/aa111-2-4}{10.4064/aa111-2-4}.

\bibitem{taoteravainen2025} T.~Tao and J.~Ter\"av\"ainen,
  \href{https://arxiv.org/abs/2512.01739v2}{\emph{Quantitative correlations and
  some problems on prime factors of consecutive integers}}, arXiv:2512.01739v2
  (submitted December 2025, revised April 2026).

\bibitem{fan2026comment} Steve Fan,
  \emph{Comment on Erd\H{o}s Problem \#269, thread 269, post 7218} (2026),
  \href{https://www.erdosproblems.com/forum/thread/269#post-7218}{source}. 26 June 2026.

\bibitem{mahler1929} K.~Mahler, \emph{Arithmetische Eigenschaften der L\"osungen
  einer Klasse von Funktionalgleichungen}, Math.\ Ann.\ \textbf{101} (1929),
  342--367.

\bibitem{adamczewskifaverjon2016} B.~Adamczewski and C.~Faverjon,
  \href{https://arxiv.org/abs/1508.07158v2}{\emph{M\'ethode de Mahler : relations
  lin\'eaires, transcendance et applications aux nombres automatiques}},
  arXiv:1508.07158v2 (2016), Theorem~1.1, p.~3.

\bibitem{nishioka1996} Ku.~Nishioka, \emph{Mahler Functions and Transcendence},
  Lecture Notes in Math.\ 1631, Springer, Berlin, 1996.

\bibitem{boesdarsterdos1981} D.~Boes, R.~Darst and P.~Erd\H{o}s,
  \href{https://www.tandfonline.com/doi/abs/10.1080/00029890.1981.11995266}{\emph{Fat,
  symmetric, irrational Cantor sets}}, Amer.\ Math.\ Monthly \textbf{88} (1981),
  340--341.

\bibitem{rstw2019} A.~Rahm, N.~Solomon, T.~Trauthwein and B.~Weiss,
  \href{https://arxiv.org/abs/1909.01198}{\emph{The distribution of rational
  numbers on Cantor's middle thirds set}}, arXiv:1909.01198.

\bibitem{chowvarjuyu2024} S.~Chow, P.~Varj\'u and H.~Yu,
  \href{https://arxiv.org/abs/2402.18395}{\emph{Counting rationals and
  diophantine approximation in missing-digit Cantor sets}}, arXiv:2402.18395.

\bibitem{plectis68} W.~Cook,
  \href{erdos-68-factorial-denominator-irrationality.pdf}{\emph{Two
  incomparable denominator exclusions for $\sum_{n\ge2}(n!-1)^{-1}$}},
  Erd\H{o}s Problem Notes, 2026.

\bibitem{plectis243} W.~Cook,
  \href{erdos-243-reciprocal-tail-rigidity.pdf}{\emph{Bounded
  increments and rational reciprocal sums}}, Erd\H{o}s Problem Notes, 2026.

\bibitem{plectis249} W.~Cook,
  \href{erdos-249-binary-totient-series.pdf}{\emph{Bases and
  integral relations for the $k$-kernel of Euler's totient}}, Erd\H{o}s Problem
  Notes, 2026.

\bibitem{plectis251} W.~Cook,
  \href{erdos-251-prime-gap-dyadic-series.pdf}{\emph{Sparse
  congruence-preserving perturbations of dyadic series}}, Erd\H{o}s Problem
  Notes, 2026.

\bibitem{plectis257} W.~Cook,
  \href{erdos-257-mersenne-support-subseries.pdf}{\emph{Weighted
  support criteria for reciprocal Mersenne subseries}}, Erd\H{o}s Problem Notes,
  2026.

\bibitem{plectis269} W.~Cook,
  \href{erdos-269-three-prime-running-lcm.pdf}{\emph{No
  finite separable representation at three prime generators}}, Erd\H{o}s Problem
  Notes, 2026.

\bibitem{plectis1049} W.~Cook,
  \href{erdos-1049-rational-base-lambert.pdf}{\emph{Zudilin's
  forms at rational bases and the exact normalised Hankel order}}, Erd\H{o}s
  Problem Notes, 2026.

\bibitem{plectisinvestigation} W.~Cook,
  \emph{Choices, contraction and rational membership}, exact computation with saved outputs in \texttt{research/experiments/choices\_contraction/} of the repository,
  2026.

\bibitem{plectisrecord} W.~Cook,
  \href{erdos-synthesis-reading-together-record.pdf}{\emph{Reading
  eight Erd\H{o}s problems together: a working record}}, 2026.
\end{thebibliography}

\section*{Declaration of generative AI use}
The mathematics and text of this note were produced with large language model
agents under the direction of the author, who is responsible for every claim.
Each statement above names its evidence: an ordinary proof printed here, a
result quoted from a companion note with the status given there, an exact
computation with saved outputs, or a cited external theorem.  The record
\cite{plectisrecord} lists what was checked, by which means, and what was found
to be already known.
\end{document}
